\section{Quantum Coherence Functions}

\subsection{Classical Coherence Functions}
Imagine a two slit experiment. For a source with bandwidth $\Delta \omega$ and path difference $\Delta s$, coherence will occur for $\Delta s \leq c/\Delta \omega$, with $\Delta s_\text{coh} = c/\Delta\omega$ the coherence length, and $\Delta t_\text{coh} = \Delta s_\text{coh}/c$ the coherence time. Let $K_1,K_2$ be complex geometric factors corresponding to the path lengths. The ergodic hypothesis tell us that the time average is equivalent to the ensemble average. Using $x_i = \vec{r}_i, t_i$ we can write the first order coherence function
\begin{equation}
    \coherence (x_1, x_2) = \frac{\expval{E^*(x_1)E(x_2)}}{\sqrt{\langle\abs{E(x_1)}^2\rangle\langle\abs{E(x_2)}^2}\rangle},
\end{equation}
and the intensity at the slits is defined as $I_i = \abs{K_i}^2 \langle \abs{E(x_i)}^2 \rangle$. Writing $K_j = \abs{K_j}\exp(i\psi_j)$ and $\coherence(x_1,x_2) = \abs{\coherence(x_1,x_2)}\exp(i\Phi_{12})$ we have the intensity at the detector 
\begin{equation}
    I(\vec{r}) = I_1 + I_2 +2\sqrt{I_1I_2}\abs{\coherence(x_1,x_2)}\cos(\Phi_{12} - \psi)
\end{equation}, with $\psi = \psi_1 - \psi_2 < \Delta s_\text{coh}$. Furthermore, $\abs{\coherence} = 1$ is complete coherence, $0 < \abs{\coherence} < 1$ is partial coherence, and $\abs{\coherence}=0$ is complete decoherence. Reyleigh's fringe visibility is then $\mathcal{V} = 2\sqrt{I_1I_2}\abs{\coherence}/(I_1 + I_2)$. For a source with coherence time $\tau_0$, inversly related to the spectral line width, and assuming a Lorentzian power spectrum, we obtain the autocoherence function $\coherence(\tau) = \exp(i\omega_0\tau - \abs{\tau}/\tau_0)$. \question{Why do we so often see a Lorentzian distribution in atomic physics? What is the physical explanation of it?} 

\subsection{Quantum Coherence Functions}
We consider a detector constructed from a single atom smaller than the wavelength (dipole approximation). Absorption leads to ionization and detection of the photoelectron, thus detecting the photon. The detector is coupled to the field by the dipole interaction (Sec \ref{sec:interaction_quantized}), and we assume that the excited photoelectron may be described by the wave function of a plane wave. We can define the first order correlation function as $G^{(1)}(x_1, x_2) = \tr(\density \hat{E}^-(x_1) \hat{E}^+(x_2))$, where $\density$ is the density function of the field and $\hat{E}^+$ is the positive part of the field as defined in Sec. \ref{sec:multimode}. The detector response from light arriving from slit $i$ is $G^{(1)}(x_i, x_i) =\tr(\density \hat{E}^-(x_i) \hat{E}^+(x_i))$. The normalized quantum coherence function is $\qcoherence(x_1, x_2) = G^{(1)}(x_1,x_2) / \sqrt{G^{(1)}(x_1,x_1)G^{(1)}(x_2,x_2)}$, with the same types of coherences as in the classical case. For a single mode number state, or coherent state we get that $\abs{\qcoherence} = 1$.

\subsection{Young's Interference}
If we have a laser beam incident on two pinholes close together, we will observe an interference pattern with fringes for angles $\Phi =2\pi m$ with $m$ an integer. The intensity falls of at a speed proportional to $1/r^2$ from the central fringe. If the incoming light is a number state or a coherent state, the interference pattern is the same, only the max intensity being shifted by some amount due to the average photon number. Note that the result will be different if the light incident on each pinhole comes from different sources.

\subsection{Higher-Order Coherence Functions}
First order coherence experiments cannot distinguish between states with different photon number distributions if their spectral distributions are the same. Hanbury Brown and Twiss developed a coherence experiment relying on the correlation of intensities instead of the fields as Young had done. The experiment uses a beam splitter and a time delay for one of the beams, and if the time delay, $\tau$, is shorter than the beam's coherence time, it is possible to determine the photon statistics of the beam. We find that the second order coherence function is $\ccoherence(\tau) = \langle I(t)I(t+\tau)\rangle / \langle I(t) \rangle^2$ if the fields are stationary, and more generally $\ccoherence(x_1,x_2;x_2,x_1) = \langle I(x_1)I(x_2)\rangle / \langle I(x_1) \rangle\langle I(x_2)\rangle$. For chaotic light we find that $1 \leq |\ccoherence(\tau)| \leq 2$ and the Lorentzian spectrum $\ccoherence(\tau) = 1 + \exp(-2|\tau|/\tau_0)$. When the time delay goes to zero we get twice the detection rate compared to long times, we say that the photons arrive in bunched pairs, called the photon bunching effect, or the Hanbury Brown and Twiss effect. Introducing the second order quantum correlation function $\qqcorrelation(x_1,x_2;x_2,x_1) = \tr(\density \hat{E}^-(x_1) \hat{E}^-(x_2) \hat{E}^+(x_2) \hat{E}^+(x_1))$ interpreted as the ensemble average $I(x_1)I(x_2)$ and is normal ordered, we can write the second order quantum coherence function as $\qqcoherence(x_1,x_2;x_2,x_1) = \qqcorrelation(x_1,x_2;x_2,x_1) / \qcorrelation(x_1,x_1)\qcorrelation(x_2,x_2)$. It is second order coherent if $|\qcoherence| = 1$ and $\qqcoherence=1$. For a single mode field we have $\qqcoherence(\tau) = \qqcoherence(0) = 1$ which implies the same relation for a coherent state. For a single mode thermal field we have $\qqcoherence(\tau) = \qqcoherence(0) = 2$, and for a multimode thermal field $\qqcoherence(\tau) = 1 + |\qcoherence(\tau)|$. Collision-broadened light gives the Lorentzian spectrum $\qqcoherence(\tau) = 1 + \exp(-2|\tau|/\tau_0)$ just as in the classical case. A single mode field in a number state can exhibit $\qqcoherence(0) < 1$, which is classically forbidden, and is a sub-Poissonian state. We can characterize states with Mandel's Q-parameter
\begin{equation}
    Q = \frac{\expval{(\Delta \hat{n})^2} - \expval{\hat{n}}}{\expval{\hat{n}}},
\end{equation}
for which states with $Q=0$ are Poissonian, $Q>0$ are super-Poissonian, and $-1<Q<0$ are sub-Poissonian. \question{What properties does these types of states share? If two states are constructed in different ways, but both are sub-Poissonian for example, what properties do they share?}